\documentclass[10pt,a4paper]{amsart}
\usepackage[margin=19mm]{geometry}
\usepackage{amsmath,amssymb,mathtools}
\usepackage[scr=rsfs,cal=euler]{mathalfa}
\usepackage{xfrac}
\usepackage[unicode,hidelinks,bookmarks=false]{hyperref}
\newcommand{\ie}{i.e.\ }
\newcommand{\cf}{cf.\ }
\newcommand{\resp}{respectively\ }
\newcommand{\Subsection}{\subsection}
\providecommand{\nobreakdash}{\nobreak\hbox{-}}
\setlength{\emergencystretch}{2em}
\multlinegap0pt
\usepackage{bm}
\usepackage{bbm}
\usepackage{dsfont}
\usepackage{xspace}
\usepackage{cancel}
\usepackage{mathtools}
\usepackage{amsthm}
\makeatletter
\@ifundefined{Thm}{\newtheorem{Thm}{Theorem}}{}
\@ifundefined{Rem}{\newtheorem{Rem}[Thm]{Remark}}{}
\makeatother
\title[$\mathcal{L}$-Lie Algebroids over Topological Ringed Spaces]{$\mathcal{L}$-Lie Algebroids over Topological Ringed Spaces}
\author{Mainak Poddar, Abhishek Sarkar}
\date{Source: arXiv:2412.12935v3 (CC BY 4.0)}
\begin{document}
\maketitle

\noindent\emph{Source and licence.} $\mathcal{L}$-Lie Algebroids over Topological Ringed Spaces. Version arXiv:2412.12935v3 (CC BY 4.0), distributed under the Creative Commons Attribution 4.0 International licence, as recorded in the arXiv licence metadata for this exact version and shown on the abstract page https://arxiv.org/abs/2412.12935v3. Full rights notice: licenses/arxiv-2412.12935-CC-BY-4.0.txt.\par\smallskip

\noindent\textit{Editorial scope.} Counted unit: the Thm whose source label is \texttt{L-Lie alg \& dga} in the article (source0.tex lines 675--678, source label \texttt{L-Lie alg \& dga}), delivered together with the article's own complete original proof of it (source0.tex lines 679--751, one \texttt{proof} environment of 73 lines). No mathematics is written, repaired or re-derived: statement and proof are the article's own text line for line. Required context retained uncounted: loc free f.r. (lines 201--208); C-E-D differential (lines 334--338); dga-reflexive (lines 752--758). Every definition, display and label of those lines is retained unchanged. Omitted from this package and not redistributed: 1--200, 209--333, 339--674, 759--1711. Nothing inside a retained excerpt is omitted or altered: every retained body is delivered exactly as the article's lines are written, so no omission receipt of any kind accompanies this package. Citations: the retained excerpts carry no citation command at all, so no bibliography entry of the article is transcribed and no reference is redistributed. Reference adaptation: none was needed and none was made; every reference inside the delivered unit resolves inside it, so no unresolved reference or citation is produced. Printed slips kept literal: no printed word, symbol, hypothesis or number of the counted statement or of its proof was altered. Layout and class: the article's own class is \texttt{sigma} with packages including ifpdf, xy, tikz-cd; the wrapper is the lane's pinned \texttt{amsart} editorial wrapper at 10pt on A4, which restates the theorem-like environments the retained text needs and adds the article's own macro definitions that the retained text needs (none), with extra wrapper packages (bm, bbm, dsfont, xspace, cancel, mathtools). The delivered numbering is the unit's own. Assets: no retained span contains a figure environment, a graphics-inclusion command, a PDF-page inclusion, a table or a TikZ picture; no image file is copied into this package, the package declares no asset dependency and no image file is redistributed with it. Licence: version arXiv:2412.12935v3 of this article is distributed under the Creative Commons Attribution 4.0 International licence, as recorded in the arXiv licence metadata for this exact version and shown on the abstract page https://arxiv.org/abs/2412.12935v3. A full rights notice is delivered at licenses/arxiv-2412.12935-CC-BY-4.0.txt. Characters: every retained excerpt is pure ASCII, so the delivered engine, XeLaTeX with its default Latin Modern text font, reports no missing character.
\begin{Rem} \label{loc free f.r.}
	For a locally free $\mathcal{O}_X$-module $\mathcal{E}$ of finite rank $n$, the following properties hold:%
	\begin{enumerate}\itemsep=0pt
		\item For each $x \in X$, there exists an open neighborhood $U_x$ of $x$ in $X$ such that \smash{$\mathcal{E}|_{U_x} \cong \mathcal{O}_X|_{U_x}^{\oplus n}$}.
		\item For each $x \in X$, the stalk $\mathcal{E}_x$ is a free $\mathcal{O}_{X,x}$-module of rank $n$.
	\end{enumerate}
	Moreover, if $\mathcal{E}$ is equipped with a Lie algebroid structure, this induces a Lie--Rinehart algebra structure on the space of sections of $\mathcal{E}$ as well as on its stalks, retaining the above properties.\looseness=1
\end{Rem}

\begin{align}
&{\rm d}_{\mathcal{L}(U)}(\omega)(D_1\wedge \cdots \wedge D_{k+1}) \nonumber\\
& \qquad = \sum^{k+1}_{i=1}(-1)^{i+1}\mathfrak{a}(U)(D_i)(\omega(D_1\wedge \cdots \wedge \hat{D_i} \wedge \cdots \wedge D_{k+1}))\nonumber\\
&\qquad\quad + \sum_{i< j}(-1)^{i+j}\omega([D_i, D_j]\wedge D_1\wedge \cdots \wedge \hat{D_i}\wedge \cdots \wedge \hat{D_j}\wedge \cdots \wedge D_{k+1}),\label{C-E-D differential}
\end{align}

\begin{Rem} \label{dga-reflexive}
	The condition of locally free $\mathcal{O}_X$-modules of finite rank in the above result can be relaxed by considering both $\mathcal{O}_X$-modules $\mathcal{E}$ and $\mathcal{L}$ as reflexive sheaves. Under this assumption, in the second part of the proof, the restriction of the DGA morphism
\[
\psi\colon \ \bigl(\mathcal{H}{\rm om}_{\mathcal{O}_X}\bigl(\wedge^\bullet_{\mathcal{O}_X}\mathcal{L}, \mathcal{O}_X\bigr), {\rm d}_{\mathcal{L}}\bigr) \rightarrow \bigl(\mathcal{H}{\rm om}_{\mathcal{O}_X}\bigl(\wedge^\bullet_{\mathcal{O}_X}\mathcal{E}, \mathcal{O}_X\bigr), {\rm d}_{\mathcal{E}}\bigr)
\]
 to the grading $1$ induces the same $\mathcal{O}_X$-module homomorphism $\phi\colon \mathcal{E} \rightarrow \mathcal{L}$, leveraging the reflexivity property. For the first part of the proof, we do not even need the reflexive condition.
\end{Rem}

\begin{Thm} \label{L-Lie alg & dga}
	Let $(\mathcal{L}, [\cdot, \cdot]_{\mathcal{L}}, \mathfrak{a}_{\mathcal{L}})$ $($or simply $\mathcal{L})$ be a locally free Lie algebroid of finite rank over $(X,\mathcal{O}_X)$.
	A~locally free $\mathcal{O}_X$-module $\mathcal{E}$ of finite rank has an $\mathcal{L}$-Lie algebroid structure if and only if $\Omega^{\bullet}_{\mathcal{E}}:=(\wedge^{\bullet}_{\mathcal{O}_X}\mathcal{E}^*,{\rm d}_{\mathcal{E}})$ has a $\Omega^{\bullet}_{\mathcal{L}}$-differential graded algebra structure. $($For more generalities, see Remark~{\rm \ref{dga-reflexive})}.
\end{Thm}

\begin{proof}
	Let the $\mathcal{O}_X$-module $\mathcal{E}$ has an $\mathcal{L}$-Lie algebroid structure given by a Lie algebroid homomorphism
\[
\phi\colon \ (\mathcal{E}, [\cdot, \cdot]_{\mathcal{E}}, \mathfrak{a}_{\mathcal{E}}) \rightarrow (\mathcal{L}, [\cdot, \cdot]_{\mathcal{L}}, \mathfrak{a}_{\mathcal{L}}).
\]
	We consider the Chevalley--Eilenberg--de Rham complex $\Omega^{\bullet}_{\mathcal{E}}:=\bigl(\wedge^{\bullet}_{\mathcal{O}_X}\mathcal{E}^*,{\rm d}_{\mathcal{E}}\bigr)$ associated to the Lie algebroid $(\mathcal{E}, [\cdot, \cdot]_{\mathcal{E}}, \mathfrak{a}_{\mathcal{E}})$ $($or simply $\mathcal{E})$. Note that, the anchor $\mathfrak{a}_{\mathcal{E}}=\mathfrak{a}_{\mathcal{L}} \circ \phi$ and the differential~${\rm d}_{\mathcal{E}}$ is given by the composition of the maps
\[
\mathcal{O}_X \overset{{\rm d}_X}{\longrightarrow}\Omega^1_X \overset{\mathfrak{a}_{\mathcal{L}}^*}{\longrightarrow} \Omega^1_{\mathcal{L}}:= \mathcal{L}^* \overset{\phi^*}{\longrightarrow} \mathcal{E}^*=:\Omega^1_{\mathcal{E}},
\]
 i.e., ${\rm d}_{\mathcal{E}}$ is the canonical extension of the map $\phi^* \circ \mathfrak{a}_{\mathcal{L}}^* \circ {\rm d}_X$, where
\[
{\rm d}_{\mathcal{E}}f (D)={\rm d}_Xf(\mathfrak{a}_{\mathcal{L}}(\phi(D)))= \mathfrak{a}_{\mathcal{L}}(\phi(D))(f)\qquad \text{for all $D \in \mathcal{E}$ and $f \in \mathcal{O}_X$}.
\]
We show that the induced map forms a homomorphism of the (sheaves of) cochain complexes
\[
\widetilde{\phi^*}\colon \ \Omega^{\bullet}_{\mathcal{L}} \rightarrow \Omega^{\bullet}_{\mathcal{E}}
\] is defined by
	$\widetilde{\phi^*}(\omega_1 \wedge \cdots \wedge \omega_k)= (\omega_1 \circ \phi )\wedge \cdots \wedge (\omega_k \circ \phi )$,
	for $\omega_1, \dots, \omega_k \in \mathcal{L}^*$ as follows. Consider $\tilde{\phi}:=\wedge^k \phi$, thus
	for a section $\omega$ of $ \Omega^{k}_{\mathcal{L}}:= \wedge^k_{\mathcal{O}_X} \Omega^{1}_{\mathcal{L}}$ and sections $s_1, \dots, s_{k+1}$ of $\mathcal{E}$ we get the following identities	(see (\ref{C-E-D differential}))	
	\begin{align*}
			& 	{\rm d}_{\mathcal{E}}\bigl(\widetilde{\phi^*}(\omega)\bigr)(s_1 \wedge \cdots \wedge s_{k+1})\\
			&\qquad = {\rm d}_{\mathcal{E}}\bigl(\omega \circ \tilde{\phi}\bigr)(s_1 \wedge \cdots \wedge s_{k+1})\\
			&\qquad = \sum^{k+1}_{i=1}(-1)^{i+1}\mathfrak{a}_{\mathcal{E}}(s_i)\bigl(\bigl(\omega \circ \tilde{\phi}\bigr) (s_1\wedge \cdots \wedge \widehat{s_i} \wedge \cdots \wedge s_{k+1})\bigr)\\
			&\qquad\quad
			+ \sum_{i< j}(-1)^{i+j} \bigl(\omega \circ \tilde{\phi}\bigr)([s_i, s_j]_{\mathcal{E}} \wedge s_1\wedge \cdots \wedge \widehat{s_i}\wedge \cdots \wedge \widehat{s_j}\wedge \cdots \wedge s_{k+1})\\
			&\qquad = \sum^{k+1}_{i=1}(-1)^{i+1}\mathfrak{a}_{\mathcal{L}}(\phi(s_i))\bigl(\omega \bigl(\phi(s_1) \wedge \cdots \wedge \widehat{\phi(s_i)} \wedge \cdots \wedge \phi(s_{k+1})\bigr)\bigr)\\
			&
			\qquad\quad + \sum_{i< j}(-1)^{i+j} \omega \bigl([\phi(s_i), \phi(s_j)]_{\mathcal{L}} \wedge \phi(s_1) \wedge \cdots \wedge \widehat{\phi(s_i)} \wedge \cdots \wedge \widehat{\phi(s_j)}\wedge \cdots \wedge \phi(s_{k+1})\bigr)\\
			&\qquad= {\rm d}_{\mathcal{L}}(\omega) (\phi(s_1) \wedge \cdots \wedge \phi(s_{k+1}))
			 = \widetilde{\phi^*}({\rm d}_{\mathcal{L}}(\omega)) (s_1 \wedge \cdots \wedge s_{k+1})
	\end{align*}
	(this is an extension of the identity $\phi^*({\rm d}_{\mathcal{L}}(f))(s)={\rm d}_{\mathcal{L}}(f)(\phi(s))$ where $f \in \mathcal{O}_X$ and $s \in \mathcal{E}$),
	i.e., the required compatibility condition ${\rm d}_{\mathcal{E}} \circ \widetilde{\phi^*}= \widetilde{\phi^*} \circ {\rm d}_{\mathcal{L}}$ holds. By the construction, the map $\widetilde{\phi^*}$ is also a graded algebra map.
	Thus, it provides a $\Omega^{\bullet}_{\mathcal{L}}$-differential graded algebra structure on $\Omega^{\bullet}_{\mathcal{E}}$.
	
Suppose, $\Omega^{\bullet}_{\mathcal{E}}:=\bigl(\wedge^{\bullet}_{\mathcal{O}_X}\mathcal{E}^*,{\rm d}_{\mathcal{E}}\bigr)$ is a $\Omega^{\bullet}_{\mathcal{L}}$-differential graded algebra with connecting morphism $\psi\colon \Omega^{\bullet}_{\mathcal{L}} \rightarrow \Omega^{\bullet}_{\mathcal{E}}$.
	Consider the restriction map $ \psi|_{\mathcal{L}^*}\colon \mathcal{L}^* \rightarrow \mathcal{E}^*$ and dualizing it we get a homomorphism $\phi\colon \mathcal{E}\rightarrow \mathcal{L}$ of $\mathcal{O}_X$-modules, since both $\mathcal{E}$ and $\mathcal{L}$ are locally free $\mathcal{O}_X$-modules of finite rank (see Remark~\ref{loc free f.r.}).
	Define the $\mathcal{O}_X$-module homomorphism
	$\mathfrak{a}_{\mathcal{E}}\colon \mathcal{E}\rightarrow \mathcal{T}_X$ as
\[
\mathfrak{a}_{\mathcal{E}}(D)(f):={\rm d}_{\mathcal{E}}f (D)
\]
	(well definedness of the map is given by the derivation property of the differential ${\rm d}_{\mathcal{E}}$).
	Since ${\rm d}_{\mathcal{E}}=\psi|_{\mathcal{L}^*} \circ \mathfrak{a}_{\mathcal{L}}^* \circ {\rm d}_X= (\mathfrak{a}_{\mathcal{L}}\circ \phi)^* \circ {\rm d}_X$ and thus $\mathfrak{a}_{\mathcal{E}}(D)(f)= (\mathfrak{a}_{\mathcal{L}}\circ \phi)^* ({\rm d}_Xf)(D)={\rm d}_Xf((\mathfrak{a}_{\mathcal{L}}\circ \phi)(D))=((\mathfrak{a}_{\mathcal{L}}\circ \phi)(D))(f)$ for all $f \in \mathcal{O}_X$, $D\in \mathcal{E}$. Hence, the map
	$\mathfrak{a}_{\mathcal{E}}$ should split as $\mathfrak{a}_{\mathcal{L}}\circ \phi$.
	Consider any two sections $D_1, D_2 \in \mathcal{E}$ and a section $\omega \in \mathcal{E}^*$, the Lie bracket $[\cdot,\cdot]_{\mathcal{E}}$ on $\mathcal{E}$ is given by the following identity:
\[
\omega([D_1, D_2]_{\mathcal{E}}):= \mathfrak{a}_{\mathcal{E}}(D_1)(\omega(D_2))- \mathfrak{a}_{\mathcal{E}}(D_2)(\omega(D_1))- {\rm d}_{\mathcal{E}}\omega (D_1 \wedge D_2).
\]
	The map $\mathfrak{a}_{\mathcal{E}}$ is a $\mathbb{K}_X$-Lie algebra homomorphism as follows:
\begin{align*}
\mathfrak{a}_{\mathcal{E}}([D_1,D_2]_{\mathcal{E}})(f)&={\rm d}_{\mathcal{E}}f([D_1,D_2]_{\mathcal{E}}) =\mathfrak{a}_{\mathcal{E}}(D_1)({\rm d}_{\mathcal{E}}f(D_2))-\mathfrak{a}_{\mathcal{E}}(D_2)({\rm d}_{\mathcal{E}}f(D_1))\\
& =\mathfrak{a}_{\mathcal{E}}(D_1)(\mathfrak{a}_{\mathcal{E}}(D_2)(f)) -\mathfrak{a}_{\mathcal{E}}(D_2)(\mathfrak{a}_{\mathcal{E}}(D_1)(f)) =[\mathfrak{a}_{\mathcal{E}}(D_1),\mathfrak{a}_{\mathcal{E}}(D_2)]_c(f)
\end{align*}
	for all $D_1,D_2 \in \mathcal{E}$ and for all $f \in \mathcal{O}_X$.
	To show the Leibniz identity for the map $\mathfrak{a}_{\mathcal{E}}$, we prove that for $f \in \mathcal{O}_X$ and $D_1,D_2 \in \mathcal{E}$ the following identity holds for all $\omega \in \mathcal{E}^*$:
	\begin{align} \label{Leiniz dga}
		\omega([D_1, f D_2]_{\mathcal{E}})= \omega(f [D_1,D_2]_{\mathcal{E}}+\mathfrak{a}_{\mathcal{E}}(D_1)(f) D_2).
	\end{align}
	Notice that, $(\mathcal{E}, [\cdot, \cdot]_{\mathcal{E}}, \mathfrak{a}_{\mathcal{E}})$ forms an $\mathcal{L}$-Lie algebroid, using the following identities (implies identity~(\ref{Leiniz dga})):
	\begin{align*}
			\omega([D_1, f D_2]_{\mathcal{E}})
			&= \mathfrak{a}_{\mathcal{E}}(D_1)(f \omega(D_2))- f \mathfrak{a}_{\mathcal{E}}(D_2)(\omega(D_1))-f {\rm d}_{\mathcal{E}}\omega (D_1 \wedge D_2)\\
			&=
			\mathfrak{a}_{\mathcal{E}}(D_1)(f) \omega(D_2)\hspace{-0.8pt}+\hspace{-0.8pt} f \mathfrak{a}_{\mathcal{E}}(D_1)(\omega(D_2)))\hspace{-0.8pt}-\hspace{-0.8pt} f \mathfrak{a}_{\mathcal{E}}(D_2)(\omega(D_1))\hspace{-0.8pt}-\hspace{-0.8pt}f {\rm d}_{\mathcal{E}}\omega (D_1 \hspace{-0.5pt}\wedge\hspace{-0.5pt} D_2)\\
			&=
			f (\mathfrak{a}_{\mathcal{E}}(D_1)(\omega(D_2))- \mathfrak{a}_{\mathcal{E}}(D_2)(\omega(D_1))-{\rm d}_{\mathcal{E}}\omega (D_1 \wedge D_2))+\mathfrak{a}_{\mathcal{E}}(D_1)(f) \omega(D_2)\\
			&=
			f \omega([D_1, D_2]_{\mathcal{E}}) + \mathfrak{a}_{\mathcal{E}}(D_1)(f) \omega(D_2)
			=\omega(f [D_1,D_2]_{\mathcal{E}}+\mathfrak{a}_{\mathcal{E}}(D_1)(f) D_2).\tag*{\qed}
	\end{align*}\renewcommand{\qed}{}	
\end{proof}

\end{document}
